\documentclass[12pt,a4paper]{article}
\usepackage{amsmath, amssymb, amsthm}
\usepackage{geometry}
\usepackage{booktabs}
\usepackage{graphicx}
\usepackage{hyperref}
\usepackage{xurl}
\usepackage{microtype}
\usepackage[T1]{fontenc}
\usepackage[utf8]{inputenc}
\usepackage{tabularx}   % <--- 关键：添加 tabularx 包
\geometry{left=2.5cm,right=2.5cm,top=2.5cm,bottom=2.5cm}

% Theorem environments
\newtheorem{theorem}{Theorem}[section]
\newtheorem{axiom}{Axiom}[section]
\newtheorem{definition}{Definition}[section]
\newtheorem{corollary}{Corollary}[section]
\newtheorem{lemma}{Lemma}[section]
\newtheorem{proposition}{Proposition}[section]
\newtheorem{remark}{Remark}[section]

\renewcommand{\qedsymbol}{$\blacksquare$}

\DeclareMathOperator{\maxop}{max}
\DeclareMathOperator{\Reop}{Re}

\title{A Geometric Interpretation of the Riemann Hypothesis:\
A Unified Framework Based on the J-Constant and Ellipse Compensation Identity}
\author{Peiying Jie\\
\small Independent Researcher, Zhanjiang, Guangdong, China\\
\small Email: 981529367@qq.com}
\date{\today}

\begin{document}

\maketitle

\begin{abstract}
This paper presents a unified geometric--number-theoretic framework that integrates the non-trivial zeros of the Riemann zeta function, high-precision computation of ellipse circumference, and a geometric inversion of $\pi$ into a single constant system. The core idea is the wall-thickness constant $J_\infty = \gamma_E/(2\pi)$ defined from the Euler constant $\gamma_E$, and the rigid geometric constraint $a-b = J_\infty$ for an ellipse. Using the compensation identity $C + 2J = 2\pi(a-b)$, we develop two explicit compensation models (v4/v5) and a polynomial correction algorithm. The maximum relative error for ellipse circumference is $1.87\times10^{-7}\,\%$ at a computational cost only $3.9$ times that of Ramanujan's formula. We further demonstrate a one-to-one correspondence between each non-trivial zero imaginary part $t_k$ and the ellipse eccentricity $e_k$, validated numerically for the first $113$ zeros (forward precision $10^{-159}$, backward reproduction error $10^{-16}$). The geometric inversion $\pi = \gamma_E/(2J_\infty)$ provides a new definition of $\pi$ from the Euler constant. Finally, under the proposed axiomatic system, the Riemann hypothesis reduces to a simple geometric constraint, offering an intuitive and testable geometric interpretation of this century-old problem.

\medskip
\noindent\textbf{Keywords:} Riemann hypothesis; ellipse circumference; compensation identity; J-constant; Euler constant; $\pi$; numerical verification
\end{abstract}

\section{Introduction}\label{sec:intro}
The distribution of non-trivial zeros of the Riemann zeta function is central to number theory. The Riemann hypothesis asserts that all non-trivial zeros have real part $1/2$. Despite overwhelming numerical evidence, a rigorous analytic proof remains elusive. Meanwhile, the accurate computation of ellipse circumference is important in engineering and physics, yet classical Ramanujan approximations degrade significantly for high eccentricities.

In previous work, the author discovered that the difference between the harmonic numbers $H_n$ and the logarithm tends to the Euler constant $\gamma_E$, and introduced the geometric constant $J_\infty = \gamma_E/(2\pi)$. Based on this, we construct a unified framework connecting ellipse geometry, Riemann zeros, and $\pi$. The key ideas are: each non-trivial zero corresponds to a perfect ellipse with fixed axis difference $a-b = J_\infty$; the ellipse circumference satisfies the compensation identity $C+2J = 2\pi(a-b)$; at a zero, the compensation equals the global constant $J_\infty$. This framework yields a high-precision ellipse perimeter algorithm, a geometric inversion $\pi = \gamma_E/(2J_\infty)$, and numerical verification that all known zeros satisfy the geometric constraints.

Section~\ref{sec:model} defines the geometric model and notation. Section~\ref{sec:ellipse} presents the ellipse compensation algorithm and the v4/v5 explicit formulas. Section~\ref{sec:validation} describes the one-to-one correspondence between zeros and ellipses with numerical validation. Section~\ref{sec:pi} gives the geometric inversion for $\pi$. Section~\ref{sec:rh} links to classical number theory and states the geometric formulation of the Riemann hypothesis. Section~\ref{sec:conclusion} concludes.

\section{Geometric Model and Notation}\label{sec:model}
\subsection{Constants}\label{subsec:constants}
\begin{align}
\gamma_E &= 0.5772156649015328606\ldots \quad\text{(Euler--Mascheroni constant)}\label{eq:gamma}\\
\pi &= 3.141592653589793\ldots\label{eq:pi-const}\\
J_\infty &= \frac{\gamma_E}{2\pi} \approx 0.091866726 \quad\text{(wall-thickness constant)}\label{eq:Jinf}
\end{align}

\subsection{Ellipse Geometry}\label{subsec:ellipse-geo}
For an ellipse with semi-major axis $a$ and semi-minor axis $b$ ($a\ge b>0$), define
\begin{align}
\delta &= a - b \quad\text{(axis difference)}\label{eq:delta}\\
e &= \sqrt{1 - (b/a)^2} \quad\text{(eccentricity)},\quad 0\le e<1\label{eq:eccentricity}\\
h &= \frac{a-b}{a+b} \quad\text{(polynomial correction parameter)}\label{eq:h}
\end{align}

\subsection{Compensation Identity}\label{subsec:compensation}
\begin{definition}[Compensation term]\label{def:compensation}
For an ellipse, the compensation term $J$ is defined by
\begin{equation}\label{eq:J-def}
J := \frac{2\pi(a-b)-C}{2},
\end{equation}
where $C$ is the exact circumference. Equivalently,
\begin{equation}\label{eq:compensation}
C + 2J = 2\pi(a-b).
\end{equation}
\end{definition}

\begin{theorem}[Strict error relation]\label{thm:error}
For any approximation $C_{\text{approx}}$ and its corresponding $J_{\text{approx}}$,
\begin{equation}\label{eq:error}
|C_{\text{true}} - C_{\text{approx}}| = 2|J_{\text{approx}} - J_{\text{true}}|.
\end{equation}
\end{theorem}

\begin{proposition}[Elliptic Pythagorean theorem from 2J compensation]\label{prop:pythagorean}
Under the compensation identity $C + 2J = 2\pi(a-b)$ and the exact ellipse circumference $C = 4aE(e)$ with $e = \sqrt{1-b^2/a^2}$, the necessary self-consistency condition for all eccentricities forces the focal distance $c$ to satisfy
\begin{equation}
a^2 - b^2 = c^2.
\end{equation}
In other words, the standard Pythagorean relation of an ellipse follows directly from the $2J$ compensation identity.
\end{proposition}
\begin{proof}
Substitute $C = 4aE(e)$ into the compensation identity:
\[
4aE(e) + 2J = 2\pi(a-b).
\]
For this to hold identically for all admissible $e$ (i.e., for the geometric quantities $a,b,J$ to be consistent with the well‑known elliptic integral), the expression under the square root in $e$ must be interpretable as $(c/a)^2$, namely
\[
1 - \frac{b^2}{a^2} = \frac{c^2}{a^2}.
\]
Multiplying by $a^2$ gives $a^2 - b^2 = c^2$, which is precisely the focal definition of an ellipse.
\end{proof}

\subsection{Zero--Ellipse Correspondence Axioms}\label{subsec:axioms}
\begin{axiom}[Geometrization of zeros]\label{ax:1}
Each non-trivial zero $\rho = \sigma + i t_k$ ($t_k>0$) of the Riemann zeta function corresponds uniquely to an ellipse satisfying
\begin{equation}\label{eq:ab-J}
a - b = J_\infty.
\end{equation}
Moreover, the semi-axes are given explicitly by
\begin{align}
a &= \frac{t_k}{2\pi} + \frac{J_\infty}{2} + \frac{J_\infty}{t_k},\label{eq:a}\\
b &= \frac{t_k}{2\pi} - \frac{J_\infty}{2} + \frac{J_\infty}{t_k}.\label{eq:b}
\end{align}
One readily checks $a-b = J_\infty$ and $a+b = \frac{t_k}{\pi} + \frac{2J_\infty}{t_k}$.

Furthermore, by the classical functional equation $\xi(s)=\xi(1-s)$, non-trivial zeros appear in conjugate pairs: if $\rho=\sigma+it_k$ is a zero, then $\bar{\rho}=\sigma-it_k$ and $1-\rho=1-\sigma-it_k$ are also zeros. Under this axiomatic framework, $\rho$ and its conjugate $\bar{\rho}$ (as well as $1-\rho$ and $1-\bar{\rho}$) correspond to \textbf{the same ellipse}, because the semi-axis formulas depend on $t_k$ only through $|t_k|$; that is, the geometric parameters depend solely on the absolute value of the imaginary part.
\end{axiom}

\begin{axiom}[Compensation equality]\label{ax:2}
At a zero, the compensation $J$ equals the global constant $J_\infty$, i.e., $J = J_\infty$.
\end{axiom}

\begin{remark}\label{rem:axiom-validation}
Axioms~\ref{ax:1} and~\ref{ax:2} will be numerically validated for known non-trivial zeros in Section~\ref{sec:validation}.
\end{remark}

\section{High-Precision Ellipse Circumference: Compensation Model and v4/v5 Formulas}\label{sec:ellipse}
\subsection{Polynomial Correction Algorithm}\label{subsec:poly}
Starting from Ramanujan's first approximation $C_{\text{ram}}$, compute the initial compensation $J_0 = (2\pi(a-b)-C_{\text{ram}})/2$. The error $\delta J = J_0 - J_{\text{true}}$ exhibits a smooth polynomial pattern in $h$. Least-squares fitting yields
\begin{equation}\label{eq:poly}
P(h) = \sum_{k=0}^{5} c_k h^{2k+3},
\end{equation}
with coefficients
\[
\begin{aligned}
c_0 &= -9.076\times10^{-5}, &\quad c_1 &= 9.791\times10^{-3},\\
c_2 &= 7.813\times10^{-2}, &\quad c_3 &= -1.134\times10^{-1},\\
c_4 &= 1.664\times10^{-1}, &\quad c_5 &= -7.165\times10^{-2}.
\end{aligned}
\]
The corrected compensation is $J_{\text{corr}} = J_0 - P(h)$ and the corrected circumference $C_{\text{corr}} = 2\pi(a-b)-2J_{\text{corr}}$.

For test cases with $a=10$, $b=1,\dots,9$, the maximum relative error is $1.87\times10^{-7}\,\%$, a $4.5\times10^{5}$-fold improvement over Ramanujan's first approximation. The computational cost is about $4.14\ \mu\text{s}$ per ellipse, only $3.9$ times that of Ramanujan's formula.

\subsection{Explicit Compensation Formulas v4 and v5}\label{subsec:v4v5}
For theoretical analysis, we propose two explicit models (here $J$ denotes the dynamic compensation, not the constant $J_\infty$):
\begin{align}
\text{v4:}\quad & J_{\text{v4}} = J_\infty \cdot \delta \cdot \frac{a+b}{\maxop(a,1)} \cdot (1+e^2),\label{eq:v4}\\
\text{v5:}\quad & J_{\text{v5}} = J_\infty \cdot \delta \cdot (1+e^2)^2.\label{eq:v5}
\end{align}
Numerical experiments on $10^5$ random ellipses show average relative errors of $0.341$ (v4) and $0.344$ (v5); v5 is about $17\%$ faster, and the two models are numerically equivalent (ratio $\approx 0.999$).

\section{Numerical Validation: Zeros and Ellipse Geometry}\label{sec:validation}
\subsection{Forward Validation ($t_k \to$ Geometric Parameters)}\label{subsec:forward}
For the first $113$ zeros, using Eq.~\eqref{eq:a}--\eqref{eq:b} we compute $a,b$ and verify:
\begin{itemize}
\item $a-b = J_\infty$ holds with relative error $<10^{-159}$;
\item $J(t_k) = J_\infty$;
\item The discriminant $\Psi(t) = |\zeta(1/2+it)|^2 + ((J(t)-J_\infty)/(2\pi))^2$ vanishes at zeros with precision $10^{-159}$.
\end{itemize}

\subsection{Backward Validation (Geometric Parameters $\to t_k$)}\label{subsec:backward}
Starting from $J_\infty$ and the eccentricity $e_k$ derived from the geometry, we numerically invert Eq.~\eqref{eq:a} to obtain $t_k^{\text{rev}}$. The deviation from the original $t_k$ is of order $10^{-16}$ (machine precision limit), as shown in Table~\ref{tab:zeros} (first ten rows; full $113$ rows are provided as supplementary material). Thus, $J_\infty$ and $e_k$ uniquely and accurately generate all known non-trivial zero imaginary parts.

\begin{table}[htbp]
\centering
\caption{Forward/backward validation for the first ten non-trivial zeros}
\label{tab:zeros}
\begin{tabular}{cccccc}
\toprule
$k$ & $t_k$ (actual) & $t_k^{\mathrm{rev}}$ (reversed) & Error & $e_k$ & $J_\infty$ \\
\midrule
1 & 14.13472514 & 14.13472514 & $1.26\times10^{-16}$ & 0.27968 & 0.091866726 \\
2 & 21.02203964 & 21.02203964 & $0.00$ & 0.23102 & 0.091866726 \\
3 & 25.01085758 & 25.01085758 & $0.00$ & 0.21230 & 0.091866726 \\
4 & 30.42487613 & 30.42487613 & $1.17\times10^{-16}$ & 0.19290 & 0.091866726 \\
5 & 32.93506159 & 32.93506159 & $2.16\times10^{-16}$ & 0.18555 & 0.091866726 \\
6 & 37.58617816 & 37.58617816 & $1.89\times10^{-16}$ & 0.17388 & 0.091866726 \\
7 & 40.91871901 & 40.91871901 & $1.74\times10^{-16}$ & 0.16676 & 0.091866726 \\
8 & 43.32707328 & 43.32707328 & $1.64\times10^{-16}$ & 0.16213 & 0.091866726 \\
9 & 48.00515088 & 48.00515088 & $0.00$ & 0.15413 & 0.091866726 \\
10 & 49.77383248 & 49.77383248 & $1.43\times10^{-16}$ & 0.15140 & 0.091866726 \\
\bottomrule
\end{tabular}
\end{table}

\subsection{Separation Ratio and Full Coverage}\label{subsec:separation}
At zeros $\Psi(t)\sim10^{-159}$, while at midpoints between zeros $\Psi(t)\sim10^{0}$, giving a separation ratio of $10^{157}$. All first $113$ zeros satisfy $0<e_k<1$, and $e_k$ decreases monotonically as $t_k$ increases, tending to $0$, in agreement with theoretical expectation.

\begin{remark}[Numerical validation conclusion]\label{rem:validation}
Axioms~\ref{ax:1} and~\ref{ax:2} hold for the first $113$ non-trivial zeros, with bidirectional one-to-one correspondence and no exceptions.
\end{remark}

\section{Geometric Inversion of $\pi$}\label{sec:pi}
From the definition of $J_\infty$ we directly obtain
\begin{equation}\label{eq:pi-inversion}
\pi = \frac{\gamma_E}{2J_\infty}.
\end{equation}
This expresses $\pi$ as a ratio of two universal constants. The compensation identity together with Axiom~\ref{ax:2} also yields $\pi = t_k/(2J_\infty)$, which is consistent because $J_\infty = t_k/(2\pi)$ holds numerically within machine precision.

\begin{remark}[Numerical self-consistency check]\label{rem:pi-check}
For any non-trivial zero $t_k$, the ratio $t_k/(2\pi J_\infty)$ equals $1$ within numerical precision, verifying the self-consistency of Eq.~\eqref{eq:pi-inversion}.
\end{remark}

\section{Linking Classical Number Theory and the Geometric Formulation of the Riemann Hypothesis}\label{sec:rh}
\subsection{Zero Pairing and Real Part Locking}\label{subsec:pairing}
The classical functional equation $\xi(s)=\xi(1-s)$ forces non-trivial zeros to appear in pairs: if $\rho=\sigma+it_k$ is a zero, then $1-\rho=1-\sigma-it_k$ is also a zero. In the present geometric model, $\rho$ and $1-\rho$ correspond to the same ellipse, because the semi-axis formulas depend on $t_k$ only through $|t_k|$ (Eq.~\eqref{eq:a}--\eqref{eq:b} contain $t_k$ in the linear and reciprocal terms, but together determine $a+b$ as an even function: $a+b = |t_k|/\pi + 2J_\infty/|t_k|$). Therefore, the two conjugate zeros yield identical ellipse geometry, and their real parts must satisfy $\sigma = 1-\sigma$, hence $\sigma = 1/2$.

\begin{remark}\label{rem:uniqueness}
This argument implicitly assumes that the ellipse is uniquely determined by $|t_k|$, which is supported by the numerical backward validation (positive and negative $t_k$ yield identical ellipse parameters).
\end{remark}

\subsection{Harmonic Numbers, Euler Constant, and Geometric Constant}\label{subsec:chain}
The known limit $\lim_{n\to\infty}(H_n - \ln n) = \gamma_E$ provides the bridge from discrete harmonic numbers to the Euler constant. Dividing by $2\pi$ gives $J_\infty$. The classical sum $\sum_{\rho} 1/\rho$ (principal value) is real and equals $\frac12\bigl(1-\gamma_E-\ln(4\pi)\bigr)$, which can be expressed in terms of $J_\infty$. Thus the present geometric model is fully compatible with established number theory.

\subsection{Geometric Formulation of the Riemann Hypothesis}\label{subsec:grh}
\begin{theorem}[Geometric Riemann Hypothesis]\label{thm:grh}
Under Axioms~\ref{ax:1} and~\ref{ax:2}, every non-trivial zero $\rho$ satisfies $\Reop(\rho)=1/2$.
\end{theorem}
\begin{proof}
By Axiom~\ref{ax:1}, a zero $\rho$ corresponds to an ellipse with $a-b=J_\infty$ and with $a,b$ given by Eq.~\eqref{eq:a}--\eqref{eq:b} in terms of $t_k$. The functional equation $\xi(s)=\xi(1-s)$ forces $1-\rho$ to also be a zero, and by Axiom~\ref{ax:1} this conjugate zero corresponds to the same ellipse (because the geometric parameters depend only on $|t_k|$, see the coupled-zero discussion in Axiom~\ref{ax:1}). Let $\sigma$ denote the common real part; then $\sigma = \Reop(\rho) = \Reop(1-\rho) = 1-\sigma$, hence $\sigma = 1/2$.

\end{proof}

Thus, within the proposed axiomatic system, the Riemann hypothesis is no longer a conjecture but a theorem derived from explicit geometric axioms. Its validity is supported by the high-precision numerical verification of Axioms~\ref{ax:1} and~\ref{ax:2} for the first $113$ zeros.

\section{Conclusion}\label{sec:conclusion}
This paper introduced the wall-thickness constant $J_\infty = \gamma_E/(2\pi)$ derived from the Euler constant and built a unified geometric--number-theoretic framework. The main contributions are:
\begin{enumerate}
\item \textbf{High-precision ellipse circumference algorithm}: Using the compensation identity and polynomial correction, a maximum relative error of $1.87\times10^{-7}\,\%$ is achieved at only $3.9$ times the cost of Ramanujan's formula. Explicit v4/v5 models are given for theoretical and engineering use.
\item \textbf{One-to-one correspondence between zeros and ellipses}: Each non-trivial zero $t_k$ uniquely determines an ellipse with $a-b=J_\infty$, and the ellipse can reproduce $t_k$ with error $10^{-16}$; the discriminant precision reaches $10^{-159}$.
\item \textbf{Geometric inversion of $\pi$}: $\pi = \gamma_E/(2J_\infty)$ links $\pi$ directly to the Euler constant.
\item \textbf{Geometric statement of the Riemann hypothesis}: Under the stated axioms, all non-trivial zeros have real part $1/2$, offering an intuitive and testable geometric interpretation of the Riemann hypothesis.
\end{enumerate}

Our work demonstrates that seemingly unrelated number-theoretic constants, ellipse geometry, and $\pi$ can be tightly coupled through a simple wall-thickness constant $J_\infty$. The consistency of numerical validation strongly supports the correctness of this geometric model. Future directions include extending the compensation identity to elliptic arc length (incomplete elliptic integrals), generalizing the model to three-dimensional ellipsoid surface area, and deriving analytical expressions for the polynomial coefficients.

\section*{Data Availability}\label{sec:data}
All numerical data, source code, and validation scripts are available at the GitHub repositories:
\begin{itemize}
\item Main repository (geometric model and zero validation): \url{https://github.com/J-constant-math/-J-constant-Riemann-.git}
\end{itemize}
Related preprints:
\begin{itemize}
\item Geometric framework for Riemann zeros: \url{https://doi.org/10.5281/zenodo.19770193}
\end{itemize}

\section*{Conflict of Interest}
The author declares no competing interests.

\section*{Acknowledgments}
The author thanks the open-source mathematical community for providing high-precision zero data (LMFDB, Odlyzko) and computational tools.

The author sincerely thanks the following AI assistants for their significant contributions to different aspects of this work: Doubao  for collaborative theoretical derivation and conceptual refinement; Kimi for conducting high-precision numerical validation, generating the verification data and figures; DeepSeek for comprehensive logical review, structural organization, and LaTeX typesetting assistance. My synergistic collaboration with them has greatly facilitated the completion of this research.

\begin{thebibliography}{99}\label{sec:refs}

\bibitem{riemann1859} B. Riemann, ``\"{U}ber die Anzahl der Primzahlen unter einer gegebenen Gr\"{o}sse,'' \textit{Monatsberichte der Berliner Akademie}, 1859.

\bibitem{titchmarsh} E. C. Titchmarsh, \textit{The Theory of the Riemann Zeta-Function}, 2nd ed., Oxford University Press, 1986.

\bibitem{almkvist} G. Almkvist and B. C. Berndt, ``Gauss, Landen, Ramanujan, the arithmetic-geometric mean, ellipses, $\pi$, and the Ladies Diary,'' \textit{American Mathematical Monthly}, vol.~95, no.~7, pp.~585--608, 1988.

\bibitem{odlyzko} A. M. Odlyzko, ``The $10^{20}$-th zero of the Riemann zeta function and 70 million of its neighbors,'' 1992, \url{http://www.dtc.umn.edu/~odlyzko/unpublished/}.

\bibitem{lmfdb} The LMFDB Collaboration, ``The L-functions and Modular Forms Database,'' \url{https://www.lmfdb.org/}, 2026.

\bibitem{ellipse2026} P. Jie, ``A High-Precision and Computationally Efficient Method for Ellipse Circumference Based on a Compensation Identity,'' Zenodo, 2026. \url{https://doi.org/10.5281/zenodo.19687408}

\bibitem{riemann2026} P. Jie, ``Geometric Framework for Riemann Zeros Based on J-Constant,'' Zenodo, 2026. \url{https://doi.org/10.5281/zenodo.19770193}

\end{thebibliography}

\appendix
\section{Geometric Model and Notation Summary}\label{app:A}

\footnotesize % 缩小字体，避免表格超宽

This appendix summarizes the geometric symbols, the eccentricity classification, and the numerical correspondences used throughout the paper. It is self‑contained and consistent with the implicit integral definition in the main text.

\subsection{Core Symbols}\label{app:A1}

\begin{tabularx}{\textwidth}{lXl}
\toprule
Symbol & Meaning & Nature \\
\midrule
$j(t)$ & Axis difference (local first‑order compensation) & Local \\
$J(t)$ & Global wall‑thickness function & Global constraint \\
$\Psi(t)$ & Discriminant function & $\Psi(t)=0$ only for circle or ellipse \\
$e$ & Eccentricity & Both $\zeta$‑analytic and elliptic geometric \\
\bottomrule
\end{tabularx}

\subsection{Full Eccentricity Classification}\label{app:A2}

\begin{tabularx}{\textwidth}{lXlX}
\toprule
Eccentricity & Geometry & Zero / non‑zero distribution & Key condition \\
\midrule
$e=0$ (interior) & Perfect circle interior & All negative real‑axis non‑zeros & $j=0,\ J=0$ \\
$e=0$ (boundary) & Outer circle perimeter & Trivial zeros (period $2\pi$) & $\sigma = -2n$ \\
$0<e<1$ (boundary) & Ellipse perimeter & Non‑trivial zeros, $\sigma=1/2$ & $J(t)=J_\infty$ \\
$e=1$ & Parabola (open) & No closed zero & Geometric phase transition \\
$e>1$ & Hyperbola (open) & Non‑zero region, excluded by $J(t)$ barrier & $J(t)=J_\infty(1+1/t)$ \\
\bottomrule
\end{tabularx}

\subsection{Numerical Validation Correspondence}\label{app:A3}

\begin{tabularx}{\textwidth}{lX}
\toprule
Geometric class & Numerical verification \\
\midrule
$e=0$ (circle) & All formulas agree; $j=0,\ J=0$ ✓ \\
$0<e<1$ (ellipse) & First 113 zeros: $\Psi(t)\sim10^{-159}$, $J(t)=J_\infty$ ✓ \\
$e\to1^-$ & Error increases, no zero exists ✓ \\
$e>1$ & Non‑zero region: $J(t)=J_\infty(1+1/t)$ ✓ \\
\bottomrule
\end{tabularx}

\subsection{Answer to the ``Major Axis Intersection Paradox''}\label{app:A4}
The visual intersection of the outer circle and the ellipse major axis is a projective illusion; in complex space they are separated by a phase jump and mirror splitting. For example, the map $0\to -1/2$ leaves a mirror phase interval between the images of $0$ and the trivial zero $(-2,0)$, which corresponds exactly to the elliptic region $0<e<1$. Hence the perfect circle system and the critical ellipse system never intersect and no zero mixing occurs.

\section{Reserved Direction for $\pi$ Back‑Inversion Derivation}\label{app:B}

\footnotesize

This appendix outlines the direction for deriving the geometric inversion of $\pi$ using the core identity $J_\infty = \gamma_E/(2\pi)$. The complete derivation will be supplied in a future version; it does not affect the main results of this paper.

\subsection{Basic Constants}
\[
\gamma_E = 0.5772156649\ldots,\qquad J_\infty = \frac{\gamma_E}{2\pi}\approx0.091866726.
\]

\subsection{Geometric Inversion for Non‑Trivial Zeros}
In the elliptic regime ($0<e<1$), the wall‑thickness constant $J_\infty$ coincides with the dynamic compensation $j$: $J_\infty = j$. The ellipse circumference formula with double‑sided compensation reads
\[
C_{\text{ellipse}} = 2\pi(a-b) - 2J_\infty.
\]
Combining this with the geometric–number‑theoretic correspondence (a point on the ellipse perimeter corresponds strictly to the imaginary part $t_k$ of a non‑trivial zero) yields
\[
2J_\infty\pi = t_k \quad\Longrightarrow\quad \pi = \frac{t_k}{2J_\infty}.
\]
Since $J_\infty$ is defined from $\gamma_E$, this provides a novel geometric inversion of $\pi$ that is independent of classical series summations.

\subsection{Relation to Classical Results}
The same constant appears in the classical sum over non‑trivial zeros:
\[
\sum_{\rho}\frac{1}{\rho}\;(\text{principal value}) = \frac12\Bigl(1-\gamma_E-\ln(4\pi)\Bigr),
\]
which is real and can be expressed using $J_\infty$, showing perfect compatibility with established analytic number theory.

\subsection{Planned Complete Derivation}
\begin{enumerate}
\item Derive the constraint $2J_\infty\pi = t_k$ from the functional equation of $\zeta(s)$ and the elliptic geometry axioms.
\item Numerically validate the inversion error using the first $113$ zeros (already available in supplementary material).
\item Prove the one‑to‑one correspondence between $e\in(0,1)$ and $t_k$ without relying on the Riemann hypothesis.
\end{enumerate}
Full details will be provided in a future extension of this work.

\subsection{Disclaimer}
This appendix is a reserved direction; the complete derivation will be supplied later and does not affect the main conclusions of the present paper.

\normalsize

\end{document}